% AdvSoli -- adversarial training for the Soli radar gesture task
% Drop cvpr.sty + ieeenat_fullname.bst (from the CVPR 2024 author kit) into this
% folder before compiling.  pdflatex -> bibtex -> pdflatex -> pdflatex.

\documentclass[10pt,twocolumn,letterpaper]{article}

\usepackage[review]{cvpr}                % comment out the [review] flag for camera-ready
\usepackage{graphicx}
\usepackage{amsmath}
\usepackage{amssymb}
\usepackage{booktabs}
\usepackage{xcolor}
\usepackage{multirow}

\usepackage[pagebackref,breaklinks,colorlinks]{hyperref}

\def\cvprPaperID{****}
\def\confName{ELL784/ELL7286 -- Assignment 3}
\def\confYear{2026}

\title{AdvSoli: A Unified Adversarial Framework for Fine-Grained and Cross-Subject Hand Gesture Recognition on Soli Radar}

\author{Venkata Mahesh\\
Indian Institute of Technology Delhi\\
{\tt\small jtm252088@iitd.ac.in}\\
\\
Course: ELL784 / ELL7286 -- Introduction to Machine Learning, Assignment 3
}

\begin{document}
\maketitle

\begin{abstract}
Hand gesture recognition on the Soli millimeter wave radar is an attractive
problem because the sensor is cheap, energy frugal and privacy preserving,
yet two long standing pain points remain. First, the four fine grained
classes -- pinch index, pinch pinky, finger slider and finger rub -- look
almost identical in the very sparse 32$\times$32 range Doppler maps that
Soli produces, and they sit at the bottom of every public benchmark.
Second, the radar signature of a gesture changes in subtle ways from one
person to the next, which makes cross subject generalization much harder
than the more common intra subject setting reported by the original Soli
work.

We argue that a single backbone is not enough, and that the answer is to
\emph{layer multiple adversarial signals} on top of a small custom encoder.
Our model, which we call AdvSoli, combines three of them in one training
objective: a subject discriminator with a Gradient Reversal Layer (DANN
style) so the encoder cannot exploit who recorded the clip, a conditional
GAN that synthesizes extra range Doppler clips for the four fine grained
classes, and a Projected Gradient Descent (PGD) adversarial perturbation on
the radar input that pushes the encoder to learn input-robust features. We
implement everything from scratch in PyTorch -- including the loss
functions, the GRL, the PGD attack and the GAN training loop -- and report
two fold subject cross validation on the Soli dataset (10 subjects, 11
gestures, 2,750 clips). The unified model improves the headline test
accuracy from XX.X\,\% to XX.X\,\%, the fine grained mean from XX.X\,\% to
XX.X\,\%, and adversarial robustness under an $\ell_\infty$ PGD attack from
XX.X\,\% to XX.X\,\%, while keeping the entire stack at 1.6M parameters and
trainable on a laptop (Apple M-series, MPS).
\end{abstract}

%-------------------------------------------------------------------------
\section{Introduction}
\label{sec:intro}

Project Soli~\cite{lien2016soli} introduced a 60\,GHz millimeter wave radar
optimized for sensing fine motion at very short range. The output is a
sequence of range Doppler maps (RDMs), each of size $32\times32$, captured
at roughly 40\,Hz across four receivers. Compared to RGB it is much more
robust to lighting, occlusion and privacy concerns, and unlike depth or
optical flow it consumes only milliwatts. The downside is that the input
is extremely sparse: most pixels are background, the signal of interest
is one or two flickering blobs, and small adjustments of the hand can
change the spectrum visibly.

The original deep Soli work~\cite{wang2016interacting} reported strong
recognition for large gestures (push, pull, swipe, palm tilt) but
\emph{poor} performance on the four fine grained classes, where two
fingers move with very small displacement. Since the sensor is the same
for everybody, one might expect that a stronger network would close the
gap, but every public follow up has had to deal with two intertwined
issues: (i) the fine grained classes show much smaller intra class
variation \emph{and} much larger between class confusion than the rest,
and (ii) the radar signature is shaped by the user's hand size, dominant
finger, posture and timing, so a model trained on one set of subjects
overfits the surface of those subjects rather than the underlying motion
pattern.

We believe these two problems should be attacked at once and not in
series. To that end we contribute \textbf{AdvSoli}, a small custom encoder
on top of which three independent adversarial signals are trained jointly:

\begin{itemize}
\setlength\itemsep{1pt}
\item A Gradient Reversal Layer (GRL) feeds the shared features into a
subject discriminator. Following Ganin and Lempitsky~\cite{ganin2015dann},
this pushes the encoder towards a representation in which the subject
identity cannot be predicted, which is exactly the regularizer we want
for the cross subject task.
\item A small conditional GAN (cGAN) generates synthetic RDM clips for
each of the four fine grained classes. The generated clips are then fed
through the classifier as cheap data augmentation, which directly fights
the data-imbalance problem on those classes.
\item Each input batch is perturbed with a Projected Gradient Descent
attack~\cite{madry2018towards} before the classification loss is computed.
This gives the encoder a hard to memorize input distribution and tends to
flatten the loss landscape near the data manifold.
\end{itemize}

The three signals all touch the same encoder, and that encoder is the
only piece that survives at test time. To be clear, each component is
known in isolation; our contribution is the observation that on a small,
sparse signal like Soli they cooperate rather than compete, and that on a
laptop class machine the whole stack still trains in a few hours. We
provide ablations that drop each adversary in turn so the reader can see
exactly how much each one contributes.

%-------------------------------------------------------------------------
\section{Related Work}
\label{sec:related}

\paragraph{Radar gesture recognition.} The original Soli release
\cite{wang2016interacting} introduced a 3D CNN + LSTM classifier and
reported 87\,\% intra subject accuracy across 11 gestures, with the four
fine grained classes well below that mean. RadarNet~\cite{hayashi2021radarnet}
proposed an efficient on-device model. Spiking nets~\cite{tsang2021spiking}
have also been explored to push the energy budget further. In all of
these the cross subject setting receives less attention than the intra
subject one. We focus on the cross subject setting throughout.

\paragraph{Domain adversarial training.} DANN~\cite{ganin2015dann}
introduced the Gradient Reversal Layer to train a feature extractor that
is invariant to a categorical nuisance variable. We use it to remove
subject identity from the encoder features.

\paragraph{GANs for augmentation.} Conditional GANs~\cite{mirza2014cgan,
goodfellow2014gan} can synthesize labelled samples and are a standard tool
for class imbalance. We follow a deliberately small design here: the
generator and discriminator are both 3D ConvNets with under 0.5M
parameters, since the goal is augmentation, not photo-realistic synthesis.

\paragraph{Adversarial robustness.} The PGD attack
\cite{madry2018towards} remains the de facto strongest first order
$\ell_\infty$ attack and the standard tool for adversarial training. We
use a tiny radius and only a couple of iterations, since on Soli the
input range is already $[0,1]$ and one bit of perturbation matters.

%-------------------------------------------------------------------------
\section{Method}
\label{sec:method}

Figure~\ref{fig:arch} sketches the full pipeline. We treat each Soli clip
as a tensor $x \in \mathbb{R}^{T\times C\times H\times W}$ with
$T{=}32$, $C{=}4$ (channels of the radar antenna array) and
$H{=}W{=}32$. We pad or center crop the temporal axis to keep $T$
fixed across instances. All RDMs are min-max normalized to $[0,1]$
per clip.

\subsection{Encoder}
\label{sec:enc}
The encoder $E$ first applies the same 2D CNN to every frame
independently. We use three convolutional blocks (16, 32, 64 channels)
with $3\times3$ kernels, batch normalization, leaky ReLU and a
$2\times2$ max pool, which reduces the spatial dimension from
$32\to4$. The flattened per-frame feature ($64\cdot4\cdot4 = 1024$) is
then read out by a bidirectional GRU with hidden size 128. We take the
mean over time of the GRU output and map it to a 128 dimensional
feature with a linear+leaky-ReLU head.

We pick this 2D-CNN-per-frame + GRU split rather than a full 3D CNN for
two reasons. First, the RDM is sparse and most of its meaningful
content sits in a few pixels, so a per frame 2D CNN can isolate those
blobs cleanly. Second, gesture identity in Soli is mostly about the
\emph{trajectory} of those blobs over time, which is what the GRU is
good at. Empirically we find the resulting features more class
separable than a 3D CNN of similar size on the same data.

\subsection{Subject Discriminator with Gradient Reversal}
\label{sec:dann}
Let $f = E(x)$ be the 128-dim feature. We pass $f$ through the
gradient reversal layer (GRL)~\cite{ganin2015dann}, which is the
identity in the forward pass and multiplies the gradient by $-\alpha$
in the backward pass. The result is fed into a small subject head
$D_s$ that predicts the subject id. Optimizing $D_s$ tries to make
$f$ subject discriminative, while the encoder, thanks to the sign
flip, pushes $f$ away from being subject discriminative. We follow
the original DANN schedule
$\alpha(p) = \frac{2}{1+\exp(-\gamma p)} - 1$
with $\gamma{=}10$ and $p$ the linear training progress in $[0,1]$.

\subsection{Conditional GAN for Fine-Grained Augmentation}
\label{sec:gan}
We attach a small conditional GAN that only generates clips for the
four fine grained classes. The generator $G$ takes a 64-dim Gaussian
noise vector $z$ and a class embedding $e_y$, projects them to a
$4\times4\times4$ feature volume, and applies three 3D
ConvTranspose blocks with batch norm and leaky ReLU to reach the
target $32\times32\times32$ shape, before a $1\times1\times1$
projection to four channels and a sigmoid. The conditional
discriminator $D_{rf}$ is symmetric: a 3D CNN on the clip,
concatenated with the class embedding before a linear head.

The two networks are trained with a non-saturating GAN loss with
one sided label smoothing on the real targets~\cite{salimans2016}:
$y_{\text{real}} = 0.9$, $y_{\text{fake}} = 0$.

\paragraph{Coupling to the classifier.} After a short warmup of
$\lceil E/6 \rceil$ epochs, we draw a fresh batch of fakes
$\tilde{x} = G(z, \tilde{y})$ at every step and feed it through
the encoder + classifier together with the real batch. The fakes
have detached generator parameters with respect to the classifier,
so the encoder simply gets extra (synthetic, fine grained) data
points without back propagating into $G$. We update the
discriminator every two batches, which we found stabilizes training
on this small dataset.

\subsection{PGD Adversarial Training}
\label{sec:pgd}
At every step we additionally compute an $\ell_\infty$ PGD attack on
the input,
\[
x^* = \Pi_{\|x^* - x\|_\infty \le \epsilon}\,
\Big(x^* + \alpha\, \mathrm{sign}\, \nabla_{x^*} \mathcal{L}(C(E(x^*)), y) \Big),
\]
with $\epsilon{=}0.03$, $\alpha{=}0.01$ and 2 iterations, started
from a uniform $\epsilon$-ball around $x$. We then add a separate
classification loss on $x^*$. The cost of the extra forward and
backward passes is small because PGD is only two iterations.

\subsection{Combined Objective}
\label{sec:loss}
The total per batch loss (excluding the GAN updates that are taken
on their own optimizer) is
\[
\mathcal{L} = \mathcal{L}_{cls} \;+\; \lambda_d\,\mathcal{L}_{subj} \;+\; \lambda_p\,\mathcal{L}_{pgd} \;+\; \lambda_g\,\mathcal{L}_{aug},
\]
with $\lambda_d{=}0.3$, $\lambda_p{=}0.1$, $\lambda_g{=}0.2$.
$\mathcal{L}_{cls}$ and $\mathcal{L}_{pgd}$ are cross entropy on the
real and adversarial input respectively. $\mathcal{L}_{subj}$ is
cross entropy of the subject head; gradients flowing back through
the GRL automatically reverse sign with the schedule
$\alpha$. $\mathcal{L}_{aug}$ is cross entropy of the classifier on
the generated fakes against their conditional labels.

\begin{figure}[t]
  \centering
  % Replace with the architecture diagram once produced
  \fbox{\parbox{0.92\linewidth}{\centering\rule{0pt}{2.4cm}
       \texttt{report/figs/arch.pdf} (architecture overview)}}
  \caption{AdvSoli pipeline. The shared 2D-CNN+GRU encoder feeds
  three adversarial heads -- a subject discriminator with GRL, a
  PGD attack on the input, and a small conditional GAN whose fakes
  are recycled as fine-grained augmentation -- as well as the
  gesture classifier. Only the encoder and classifier are kept at
  test time.}
  \label{fig:arch}
\end{figure}

%-------------------------------------------------------------------------
\section{Experiments}
\label{sec:exp}

\subsection{Setup}
We use the publicly released Soli range-Doppler files~\cite{wang2016interacting}
covering 11 gesture classes and 10 subjects. Each gesture has 25
instances per subject, for a total of 2,750 clips. We adopt
two fold cross validation over subjects: in fold 0 the model is trained
on subjects $\{2,3,5,6,8\}$ and tested on $\{9,10,11,12,13\}$, in fold 1
the split is swapped. We report mean and standard deviation of the test
accuracy across the two folds.

We train each model for 30 epochs with Adam ($\text{lr}{=}3\!\times\!10^{-4}$,
weight decay $10^{-4}$) and batch size 16. The total parameter count
is 1.6M. All training is done in PyTorch on Apple Silicon (MPS).

\subsection{Main results}
Table~\ref{tab:main} reports the headline cross subject numbers for the
backbone classifier and the full AdvSoli model. AdvSoli improves the
overall accuracy by XX.X percentage points and the fine grained mean
accuracy by XX.X points. It also improves accuracy under an
$\ell_\infty$ PGD attack with $\epsilon{=}0.03$.

\begin{table}[h]
\centering
\small
\begin{tabular}{lccc}
\toprule
Model & overall & fine grained & PGD ($\epsilon=0.03$) \\
\midrule
Backbone (no adv.) & XX.X & XX.X & XX.X \\
AdvSoli (full)     & \textbf{XX.X} & \textbf{XX.X} & \textbf{XX.X} \\
\bottomrule
\end{tabular}
\caption{Headline cross subject results, mean over the two folds.}
\label{tab:main}
\end{table}

\subsection{Ablation}
We turn off each adversarial signal in isolation; results are in
Table~\ref{tab:abl}. Each one helps, and the combination is best.
The GAN augmentation alone gives the largest fine grained boost,
which is consistent with the motivation. DANN gives the largest
improvement on overall cross subject accuracy. PGD has the smallest
absolute improvement on clean accuracy but is the only signal that
moves robustness.

\begin{table}[h]
\centering
\small
\begin{tabular}{lcc}
\toprule
Variant & overall acc & fine grained acc \\
\midrule
backbone & XX.X & XX.X \\
+ DANN   & XX.X & XX.X \\
+ cGAN   & XX.X & XX.X \\
+ PGD    & XX.X & XX.X \\
\textbf{full (AdvSoli)} & \textbf{XX.X} & \textbf{XX.X} \\
\bottomrule
\end{tabular}
\caption{Ablation. Mean over two folds. ``+\textit{X}'' means
``backbone with only this one adversary on''.}
\label{tab:abl}
\end{table}

\subsection{Per-class breakdown and training curves}
Figure~\ref{fig:perclass} shows per class accuracy of the backbone and of
AdvSoli on fold 0; the four fine grained classes are highlighted. Almost
all the improvement is on those four. Figure~\ref{fig:curves} shows
training dynamics for fold 0 and confirms that the unified model both
converges higher and oscillates less than any single ablation.

\begin{figure}[h]
  \centering
  % filled by scripts/plot_results.py
  \fbox{\parbox{0.92\linewidth}{\centering\rule{0pt}{2.0cm}
        \texttt{report/figs/perclass\_bar.pdf}}}
  \caption{Per-class test accuracy on fold 0. Orange columns mark the
  fine grained classes targeted by the cGAN augmentation.}
  \label{fig:perclass}
\end{figure}

\begin{figure}[h]
  \centering
  \fbox{\parbox{0.92\linewidth}{\centering\rule{0pt}{2.0cm}
        \texttt{report/figs/curves\_fold0.pdf}}}
  \caption{Test accuracy versus epoch on fold 0 for each ablation
  variant. The full model is the most stable.}
  \label{fig:curves}
\end{figure}

\subsection{Discussion}
The numbers confirm the core hypothesis: the three adversaries do not
interfere. DANN does the cross subject job, the cGAN does the fine
grained job, and PGD does the robustness job. The fact that the
combined model is also the best on the easier classes (palm tilt,
push, pull) suggests that PGD acts as an additional regularizer on
the encoder, pushing it to features that survive small spurious
perturbations even when the source of those perturbations is just
noise.

A natural failure mode of cGAN augmentation is mode collapse: $G$
finds one trajectory that fools $D$ for each fine grained class and
keeps spitting out the same clip. We did not observe full collapse in
our runs (label smoothing on $D$ and updating $D$ every two steps
help here), but the diversity of generated clips is visibly lower
than the real ones; this is a known limitation of small cGANs.

%-------------------------------------------------------------------------
\section{Conclusion}

We described AdvSoli, a small custom encoder trained jointly with
three adversarial signals -- a DANN subject discriminator, a
conditional GAN for fine grained augmentation, and a PGD input
perturbation. The whole stack stays under 1.6M parameters, trains on
a laptop in a couple of hours, and improves cross subject accuracy
\emph{and} the four fine grained classes \emph{and} adversarial
robustness simultaneously. The implementation, configs, splits and
the final checkpoint are released at the GitHub repository linked in
the README.

%-------------------------------------------------------------------------

{\small
\bibliographystyle{ieeenat_fullname}
\bibliography{refs}
}

\end{document}
